\section{Conclusion}
\label{ch:conclusion-and-future-work}
\subsection{General Observations}

Overall, the model shows encouraging learning behavior and seems to be able to process
long contexts efficiently. While its accuracy remains below that of transformer-based
models, it shows a possible advantage in terms to latency and scales well in 
long-context inference scenarios.\\ \\
%
The reproduced benchmark results did not fully match the reported performance, 
likely due to missing implementation details or differences in evaluation frameworks.
This points to a possible sensitivity of reported results to evaluation setup and 
preprocessing choices.
%
Logit lens analysis across stacked recurrent layers reveals behavior similar to transformer 
architectures, where intermediate representations progressively align with final predictions. 
Hidden state comparisons with Mistral-7B showed approximately 95\% CKA similarity in 
the final layers, suggesting that xLSTM develops internal representations closely aligned with those of transformer-based language models. Additionally, the model was able to identify 
and prioritize important tokens within a given sequence, attributing stronger updates to meaningful elements 
such as entities and key contextual markers. The input and forget gates seem to work in a coordinated way to 
preserve relevant information while selectively integrating new content into the memory state. \\ \\
%
Memory caching in a question-answering setup yields up to a 2.7$\times$ speedup in generation 
time by reusing precomputed memory states and avoiding redundant computation over long contexts 
accross 102 question answer pairs.
The speedup increases with context length, making the approach particularly beneficial 
for long-sequence inference. Notably, in our experiments, this improvement was achieved 
\textbf{without any observed degradation in accuracy}. However, performance still remains below 
a comparable transformer baseline (Mistral-7B).
%
\subsection{Limitations}
Despite the promising results, the model underperforms compared to a tranformer similar in size
across several benchmarks. In addition, experimental work was constrained 
by engineering limitations, including difficulties in deploying Triton kernels across 
different operating systems.

\subsection{Future Work}
This study represents an initial exploration of the model's capabilities. Future work should extend
evaluation to longer sequences and include benchmarks that explicitly test long-range memory and 
contextual retention. This would provide a more comprehensive assessment of the model's strengths 
and limitations in long-context reasoning settings.
\newpage